\section{Semantic Roles}


\begin{frame}{Situations}

\begin{itemize}
\item Noun phrases refer to entities (things, people, places, \ldots) 
\item Sentences refer to situations
  \begin{itemize}
  \item Situations can be constrained in various ways
  \item What is the event in question?
  \item Who participates in it?
  \item When did it happen?
  \item Is it ongoing or has it finished?
  \item Is our knowledge of it certain?
  \end{itemize}
\item The core of an event is typically represented by a verb or
  adjective 
  \begin{itemize}
  \item Verbs typically refer to actions (but can refer to states)
    \\ \eng{He \ul{stepped} swiftly forward \ldots} DANC
    \\  \eng{I \ul{know} you, you scoundrel!} DANC
  \item Adjectives typically refer to states
    \\ \eng{Your sister is \ul{dead}, then?} DANC
  \end{itemize}
\end{itemize}

\end{frame}

\begin{frame}{Semantic  Roles}
\MyLogo{The set of roles we use is from \citet[Chapter~6]{Saeed:2015}.}
In this section we talk about the relations between the participants
in a situation and the situation itself.

\begin{itemize}
\item  \txx{Semantic roles} are parts of the sentence that 
correspond to the participants in the situation 
described

\item  They classify relations between entities in a situation
\\[2ex] \begin{dependency}
\begin{deptext}[column sep=1em]
He \& seized \& the poker, \& and \& bent \& it \& into \& a curve \ldots \\ %with his huge brown hands.}
\end{deptext}
\deproot[edge unit distance=2.5ex]{2}{ROOT}
\depedge{1}{2}{AGENT}
\depedge{3}{2}{PATIENT}
\depedge{1}{5}{AGENT}
\depedge{6}{5}{PATIENT}
\depedge{7}{6}{GOAL}
\end{dependency}
 \newpage
\item  Also known as
\begin{itemize}
\item  Deep case \citep{Fillmore:1968}
\item  Thematic roles; Theta roles;   $\theta$-roles
\item  Participant Roles
\end{itemize}
\end{itemize}
 
\end{frame}

\begin{frame}{Roles link different alternations}
\begin{exe}
  \ex\eng{Kim patted Sandy}
  \ex\eng{Sandy was patted by Kim}
\end{exe}
\begin{itemize}
\item The semantic roles are different from the grammatical relations.
\item Which is the \txx{subject} and which the \txx{object} in these sentences?
\item What are the semantic roles of Kim and Sandy?
\item \txx{semantic dependencies}: An abstract representation of the meaning links word-senses to
  each other using semantic roles: different sentences may end up the same at this level
\\[2ex] \begin{dependency}
\begin{deptext}[column sep=1em]
Kim \& pat$_1$ \& Sandy \\ %with his huge brown hands.}
\end{deptext}
\deproot[edge unit distance=1.5ex]{2}{ROOT}
\depedge{1}{2}{AGENT}
\depedge{3}{2}{PATIENT}
\end{dependency}

\end{itemize}

\end{frame}

\begin{frame}{Semantic Roles}
\MyLogo{Some prose and examples from \citet{Bender:2013}}
\begin{itemize}
\item \txx{AGENT} (takes \eng{deliberately, on purpose, what did X do?})
  \begin{quote}
    A participant which the meaning of the verb specifies as doing or causing something,
    possibly intentionally. 
  \end{quote}
  \begin{itemize}
  \item The initiator, performer of controller of an action; typically volitional, typically animate
  \item Typically \textsc{subject}
  \end{itemize}
  \begin{exe}
      \ex\eng{\eng{\ul{Kim} kicked Sandy}}
      \ex\eng{\eng{\ul{The ogre} leaped into the fray}}
      \ex\eng{\eng{\ul{The student} watched the video}}
    \end{exe}
\item (\txx{ACTOR}) generalization of \txx{AGENT} that allows non-volitional, non-actor
  \\ if you use this, then \txx{AGENT} is restricted to animate, volitional participants
\item Find an example of AGENT in \Story{}\task
\newpage
\item  \txx{PATIENT} (\eng{What happened to X?})
  \begin{quote}
    A participant which the verb characterizes as having something
    happen to it, and as being affected by what happens to it.
    %Examples: objects of kill, eat, smash but not those of watch, hear
    %and love.
  \end{quote}
  \begin{itemize}
  \item The undergoer of an action
  \item  Undergoes change in state usually, both animate and 
    inanimate
  \item Typically \textsc{object}
  \end{itemize}
  \begin{exe}
    \ex\eng{\eng{Kim kicked \ul{Sandy}}}
    \ex\eng{\eng{The ogre ate  \ul{the dog}}}
    \ex\eng{$^\#$\eng{The student watched \ul{the video}}}
    \ex\eng{$^\#$\eng{I heard \ul{a sound}}}
  \end{exe}
\item Find an example of PATIENT in \Story{}\task
\newpage
\item  \txx{THEME}
  \begin{quote}
     A participant which is characterized as changing its position or condition, or as
being in a state or position. 
  \end{quote}
  \begin{itemize}
  \item  Moved, location or state is described
  \item Typically \textsc{object}
\end{itemize}
\begin{exe}
  \ex \eng{\eng{Hiromi put \ul{the book} on the shelf}}
  \ex \eng{\eng{Freddy gave you \ul{the chocolate}}}
  \ex \eng{\eng{\ul{The book} is on the shelf}}
  \ex \eng{\eng{\ul{The protagonist} died}}
  \ex *\eng{\eng{\ul{The dog} walked home}}
\end{exe}
\item Find an example of THEME in \Story{}\task
\newpage

\item  \txx{EXPERIENCER}
  \begin{quote}
    A participant who is characterized as aware of something.
  \end{quote}
  \begin{itemize}
  \item   Non-volitional, displaying awareness of action, state
  \item Typically \textsc{subject}
  \end{itemize}
  \begin{exe}
  \ex\eng{\ul{Liling} heard thunder}
  \ex\eng{\ul{Jo} felt sick}
  \ex\eng{The lecturer annoyed \ul{the students}}
\end{exe}
\item  \txx{STIMULUS}
  \begin{itemize}
  \item  Usually used in connection with EXPERIENCER
  \end{itemize}
  \begin{exe}
  \ex\eng{\ul{The lightning} scared them}
  \ex\eng{I don't like \ul{the lightning}}
\end{exe}
\item Find examples of EXPERIENCER and  STIMULUS in \Story{}\task

  \newpage
\item  \txx{BENEFICIARY}
  \begin{itemize}
  \item   For whose benefit the action was performed
  \item   Typically indexed by \lex{for} PP in English
    \\ or \textsc{object} in ditransitive verbs
  \end{itemize}
  \begin{exe}
  \ex\eng{They made \ul{me} a present}
  \ex\eng{They made a present \ul{for me}}
\end{exe}

\item  \txx{LOCATION}
  \begin{itemize}
  \item  Place
  \item Typically indexed by locative PPs (\lex{in, on, at}, \ldots) in English
  \end{itemize}
  \begin{exe}
  \ex\eng{I am living \ul{in Indonesia}}
  \ex\eng{It is \ul{on the table}}
\end{exe}
\item Find  examples of BENEFICIARY and LOCATION in \Story{}\task
\newpage  

\item  \txx{GOAL}
  \begin{itemize}
  \item  towards which something moves (lit or metaphor)
 \item  Typically indexed by \lex{to} PP in English 
   \\ or \textsc{object} in ditransitive
  \end{itemize}
  \begin{exe}
  \ex\eng{She handed the form \ul{to him}}
  \ex\eng{She handed \ul{him} her form}
  \end{exe}
\item  \txx{SOURCE}
 \begin{itemize}
 \item  from which something moves or originates
 \item  Typically indexed by \ex{from} PP in English
 \end{itemize}
  \begin{exe}
 \ex\eng{We gleaned this \ul{from the Internet}}  
\end{exe}
\item Find  examples of SOURCE and GOAL in \Story{}\task
\newpage
\item  \txx{INSTRUMENT}
  \begin{itemize}
  \item  Means by which action is performed
  \item  Can be indexed by \lex{with} PP in English
  \end{itemize}
  \begin{exe}
    \ex\eng{I ate breakfast \ul{with chopsticks}}
  \end{exe}
\item Find examples of INSTRUMENT in \Story{}\task
\item Some researchers have a separate COMITATIVE case which encodes a
  relationship of accompaniment between two participants in an event 
\end{itemize}


\end{frame}

\begin{frame}{PropBank Roles}

\begin{itemize}\addtolength{\itemsep}{-1ex}
\item An influential set of semantic roles comes from  PropBank \citep{Palmer:Gildea:Kingsbury:2005}
\item They have 6 core roles and 17 modifier roles
\item The core roles, meaning is defined per verb sense \\
  \begin{tabular}{lll}
    Role & Description & Example \\
    ARG0 & agent          & \\
    ARG1 & patient   & \\
    ARG2 & instrument, benefactive, attribute   & \\
    ARG3 & starting point, benefactive, attribute   & \\
    ARG4 & ending point   & \\
    ARGA & secondary agent & \eng{\ul{Kim}$_A$ \ull{trotted} \ul{her horse}$_0$}
      \end{tabular}
    \item \lex{seize.01} ``acquire (forcefully or stealthily)''
    \\  \begin{tabular}{ll}
        Arg0-PAG: &agent, entity acquiring something\\
        Arg1-PPT: &thing acquired \\
        Arg2-DIR: &acquired-from
      \end{tabular}
      \\[1.5ex] \eng{\ul{Sandy}$_0$ \ull{seized} \ul{the poker}$_1$ \ul{from Kim}$_2$}
    \end{itemize}


\end{frame}

\begin{frame}{The Modifier Roles}
\MyLogo{Can be used with any verb}
\begin{small}
\begin{verbatim}
COM: Comitative
LOC: Locative
DIR: Directional
GOL: Goal
MNR: Manner
TMP: Temporal
EXT: Extent
REC: Reciprocals
PRD: Secondary Predication
PRP: Purpose
CAU: Cause
DIS: Discourse
ADV: Adverbials
ADJ: Adjectival
MOD: Modal
NEG: Negation
DSP: Direct Speech
LVB: Light Verb
CXN: Construction
\end{verbatim}
  \end{small}

\end{frame}

\begin{frame}{Some Issues}
\begin{itemize}
\item  Every theory has a different set of roles
\item  It is hard to generalize: roles can be very word specific
\item  Roles are very under-specified, these are all PATIENT!
\begin{exe}
\ex\eng{The genie touched \ul{the lamp} with their nose.}
\ex\eng{The baby rubbed \ul{the lamp} with its hands.}
\ex\eng{The baby squeezed \ul{the rubber toy} with its hands.}
\ex\eng{She cracked \ul{the mirror} with a stone.}
\end{exe}
\end{itemize}



\end{frame}

\begin{frame}{Linking Grammatical Relations and Semantic Roles}
\begin{itemize}
\item    Semantic roles typically map onto grammatical 
  functions systematically
  \begin{itemize}
  \item  AGENT is usually the subject
  \item  PATIENT is usually the object
  \end{itemize}
\item  It is possible to predict how arguments are linked to 
the verb from their semantic roles, and hence their 
grammatical functions.
% \end{itemize}

\item Many verbs allow \txx{alternation}s ``syntactic variants with different roles''

\begin{exe}
  \ex \eng{Jo broke the ice with a pickaxe}  \sr{AGENT, PATIENT, INSTRUMENT}
  \ex \eng{The pickaxe broke the ice}  \sr{INSTRUMENT, PATIENT}
  \ex \eng{The ice broke}  \sr{PATIENT}
\end{exe}
\end{itemize}

\end{frame}

\begin{frame}{Other Predicates}

\begin{itemize}
\item Adjectives  (normally theme)
  \begin{exe}
  \ex \eng{John is tall} (THEME)
  \ex \eng{John is cold [to touch]} (THEME)
  \ex \eng{John is/feels cold} (EXPERIENCER)
    \\ different adjectives in e.g., Japanese：
    \\  \eng{冷たい} \jpn[cold (to touch)]{tsumetai}  vs \eng{寒い} \jpn[(feel) cold]{samui}
  \end{exe}
\item Predicative Copula (treat second NP as predicate)
  \begin{exe}
  \ex \eng{John is a boy} (THEME)
  \end{exe}
\item Identity Copula (reversible)
  \begin{exe}
  \ex \eng{Kim is my teacher} (THEME, THEME)?
  \ex \eng{My teacher is Kim} (THEME, THEME)?
  \end{exe}


\end{itemize}



\end{frame}

%\begin{frame}{Thematic Hierarchy}

% \begin{itemize}
% \item  The higher you are in the hierarchy the more likely 
% to be subject (then object, then indirect, then 
% argument PP, then adjunct PP)
% \end{itemize}
% \[  \mbox{AGENT} > 
% \left\{\begin{array}[c]{l} \mbox{GOAL/RECIPIENT} \\ \mbox{BENEFICIARY} \end{array} \right\}  >
% \left\{\begin{array}[c]{l} \mbox{THEME} \\ \mbox{PATIENT} \end{array} \right\}  >
% \mbox{INSTRUMENT} > 
% \mbox{LOCATION} \]

% \begin{itemize}
% \item  Generally true across languages
% \end{itemize}



% \end{frame}

\begin{frame}{Dowty's Proto-Arguments}
% \MyLogo{\citet{Dowty:1991}}
% \begin{itemize}
% \item  The \txx{Agent} Proto-Role 
% \begin{itemize}
% \item  Volitional
% \item  Sentient (and/or perceptive)
% \item  Causes event or change of state; 
% \item  Movement
% \end{itemize}
% \item  The \txx{Patient} Proto-Role
% \begin{itemize}
% \item  Change of state
% \item  Incremental theme (i.e. determines aspect)
% \item  Causally affected by event
% \item  Stationary (relative to movement of proto-agent).     
% \end{itemize}
% \end{itemize}

% \end{frame}

%\begin{frame}{Dowty's Argument Selection Principle}
% \begin{itemize}
% \item   when a verb takes a subject and an object
%   \begin{itemize}
%   \item  the argument with the greatest number of Proto-Agent 
%     properties will be the one selected as \textsc{subject}
%   \item  the one with the greatest number of Proto-Patient properties 
%     will be selected as \textsc{object}
%   \end{itemize}
% \item   Try: \lex{threw} --- ball, the man, the dog

% \item   Relatively predictive, but what about sentences 
% such as:
% \\  \eng{The hunger killed him}?
% \end{itemize}


%\end{frame}
